% !TEX program = xelatex
\documentclass[12pt]{article}
\usepackage[margin=28mm]{geometry}
\usepackage{hyperref}
\usepackage{amsmath, amsthm, amssymb, amsfonts}
\usepackage{tikz-cd}
\usepackage{fontspec}
\usepackage{unicode-math}
\usepackage{xeCJK}
\setmainfont{Latin Modern Roman}
\setmonofont{Latin Modern Mono}
\setCJKmainfont{Noto Serif CJK JP}
\usepackage{listings}
\lstset{basicstyle=\ttfamily\footnotesize, columns=fullflexible, breaklines=true, showstringspaces=false}
\newtheorem{theorem}{定理}
\newtheorem{lemma}{補題}
\theoremstyle{definition}
\newtheorem{definition}{定義}
\newtheorem{remark}{注意}

\title{ExpLawCO\_Nat — 日本語ノート}
\author{TeX generated by ChatGPT-5 Pro\\Lean source generated by CODEX (OpenAI)}
\date{2025-09-09}

\begin{document}
\maketitle

\begin{abstract}
A tiny "NatIso-like" record for a family of homeomorphisms natural in `X`.
We intentionally avoid `CategoryTheory.NatIso` to keep this independent of `TopCat`.
\end{abstract}

\section*{生成元の明示（Provenance)}
\begin{itemize}
  \item \textbf{TeX 生成}: ChatGPT-5 Pro
  \item \textbf{Lean ソース生成}: CODEX (OpenAI)
\end{itemize}

\section*{概説}
本ノートは、Lean ソースから導出される定義や命題を一覧化し、合わせて一般的な可換図式を掲示します。
内容は自動生成であり、厳密な注釈や証明は Lean 側の原本に依拠します。

\section*{サマリー}
Total lines: 135\\
Defs: 2\\
Lemmas: 0\\
Theorems: 0\\
Structures: 2\\
Inductives: 0

\section*{図式（例）}

\[
\begin{tikzcd}
A \arrow[r, "f"] \arrow[d, "g"'] & B \arrow[d, "h"] \\
C \arrow[r, "k"'] & D
\end{tikzcd}
\]



\[
\mathrm{Hom}(X\times A,\,B) \;\cong\; \mathrm{Hom}(X,\,B^A)
\]



\[
\begin{tikzcd}
1 \arrow[r, "0"] \arrow[d, "!"'] & \mathbb{N} \arrow[d, "u"] \\
X \arrow[r, "x_0"'] & X
\end{tikzcd}
\quad
\begin{tikzcd}
\mathbb{N} \arrow[r, "\mathrm{succ}"] \arrow[dr, swap, "u"] & \mathbb{N} \arrow[d, "u"] \\
& X
\end{tikzcd}
\]


\section*{定理・定義の一覧（先頭行）}
\begin{center}
\begin{tabular}{lll}
\textbf{種別} & \textbf{名前} & \textbf{シグネチャ} \\ \hline
structure & \texttt{NatHomeoInX} & \texttt{(Y Z : Type*) [TopologicalSpace Y] [TopologicalSpace Z] [LocallyCompactSpace Y] where} \\
def & \texttt{expLawCO\_natIso\_X} & \texttt{ (Y Z : Type*) [TopologicalSpace Y] [TopologicalSpace Z] [LocallyCompactSpace Y] : NatHomeoInX Y Z :=} \\
structure & \texttt{NatHomeoInY} & \texttt{(X Z : Type*) [TopologicalSpace X] [TopologicalSpace Z] [LocallyCompactSpace X] where} \\
def & \texttt{expLawCO\_natIso\_Y} & \texttt{ (X Z : Type*) [TopologicalSpace X] [TopologicalSpace Z] [LocallyCompactSpace X] : NatHomeoInY X Z :=} \\
\end{tabular}
\end{center}

\section*{Lean ソース（全文）}
\begin{lstlisting}
-- Natural-in-X packaging for the compact–open exponential law, in `ContinuousMap`/`Homeomorph` land.
-- No TopCat mixing; pointwise `@[simp]` lemmas; proofs via `rfl`/`ext`.

import Mathlib.Topology.CompactOpen
import MyProjects.Topology.B.ExpLawCO

noncomputable section
open scoped Topology

namespace MyProjects.Topology.B

/-- A tiny "NatIso-like" record for a family of homeomorphisms natural in `X`.
We intentionally avoid `CategoryTheory.NatIso` to keep this independent of `TopCat`. -/
structure NatHomeoInX (Y Z : Type*)
    [TopologicalSpace Y] [TopologicalSpace Z] [LocallyCompactSpace Y] where
  app : ∀ (X : Type*) [TopologicalSpace X] [LocallyCompactSpace X],
    C(X × Y, Z) ≃ₜ C(X, C(Y, Z))

variable {X X' Y Z : Type*}
variable [TopologicalSpace X] [TopologicalSpace X'] [TopologicalSpace Y] [TopologicalSpace Z]
variable [LocallyCompactSpace X] [LocallyCompactSpace X'] [LocallyCompactSpace Y]

/-- A thin alias exposing the app as `expLawCO_natIso_X.app _`. -/
def expLawCO_natIso_X
  (Y Z : Type*) [TopologicalSpace Y] [TopologicalSpace Z] [LocallyCompactSpace Y] :
  NatHomeoInX Y Z :=
{ app := fun X _ _ => expLawCOHomeo (X:=X) (Y:=Y) (Z:=Z) }

-- Convenience simp lemmas (pointwise)

@[simp] lemma expLawCO_natIso_X_apply
  [TopologicalSpace X] [LocallyCompactSpace X]
  [TopologicalSpace Y] [LocallyCompactSpace Y] [TopologicalSpace Z]
  (f : C(X × Y, Z)) (x : X) (y : Y) :
  (expLawCO_natIso_X (Y:=Y) (Z:=Z)).app X f x y = f (x, y) := rfl

@[simp] lemma expLawCO_natIso_X_symm_apply
  [TopologicalSpace X] [LocallyCompactSpace X]
  [TopologicalSpace Y] [LocallyCompactSpace Y] [TopologicalSpace Z]
  (g : C(X, C(Y, Z))) (x : X) (y : Y) :
  ((expLawCO_natIso_X (Y:=Y) (Z:=Z)).app X).symm g (x, y) = g x y := rfl

/-- Naturality in `X` (precomposition) — pointwise form, `@[simp]` friendly. -/
@[simp] lemma expLawCO_natIso_X_precomp_left_apply
  [TopologicalSpace X] [TopologicalSpace X'] [LocallyCompactSpace X] [LocallyCompactSpace X']
  [TopologicalSpace Y] [LocallyCompactSpace Y] [TopologicalSpace Z]
  (φ : C(X', X)) (f : C(X × Y, Z)) (x' : X') (y : Y) :
  ((expLawCO_natIso_X (Y:=Y) (Z:=Z)).app X'
      (f.comp (ContinuousMap.prodMk (φ.comp ContinuousMap.fst) ContinuousMap.snd))) x' y
  =
  ((expLawCO_natIso_X (Y:=Y) (Z:=Z)).app X f) (φ x') y := by
  simpa using
    (expLawCO_precomp_left_apply (X:=X) (X':=X') (Y:=Y) (Z:=Z) φ f x' y)

namespace MyProjects.Topology.B

/-! # Y-naturality and Z-postcomposition (pointwise `@[simp]`)

These are pointwise forms derived from the `_apply` lemmas in `ExpLawCO.lean`.
They keep everything inside `ContinuousMap`/`Homeomorph` with `rfl`/`ext` proofs.
-/

@[simp] lemma expLawCO_natIso_X_precomp_right_apply
  {X Y Y' Z : Type*}
  [TopologicalSpace X] [TopologicalSpace Y] [TopologicalSpace Y'] [TopologicalSpace Z]
  [LocallyCompactSpace X] [LocallyCompactSpace Y] [LocallyCompactSpace Y']
  (τ : C(Y', Y)) (f : C(X × Y, Z)) (x : X) (y' : Y') :
  ((expLawCO_natIso_X (Y:=Y') (Z:=Z)).app X
      (f.comp (ContinuousMap.prodMk ContinuousMap.fst (τ.comp ContinuousMap.snd)))) x y'
  =
  ((expLawCO_natIso_X (Y:=Y) (Z:=Z)).app X f) x (τ y') := by
  simpa using
    (expLawCO_precomp_right_apply (X:=X) (Y:=Y) (Y':=Y') (Z:=Z) τ f x y')

@[simp] lemma expLawCO_natIso_X_postcomp_apply
  {X Y Z Z' : Type*}
  [TopologicalSpace X] [TopologicalSpace Y] [TopologicalSpace Z] [TopologicalSpace Z']
  [LocallyCompactSpace X] [LocallyCompactSpace Y]
  (ψ : C(Z, Z')) (f : C(X × Y, Z)) (x : X) (y : Y) :
  ((expLawCO_natIso_X (Y:=Y) (Z:=Z')).app X (ψ.comp f)) x y
  =
  ψ (((expLawCO_natIso_X (Y:=Y) (Z:=Z)).app X f) x y) := by
  simpa using
    (expLawCO_postcomp_apply (X:=X) (Y:=Y) (Z:=Z) (Z':=Z') ψ f x y)

end MyProjects.Topology.B

-- Symmetric packaging: Y‑naturality as a tiny record

noncomputable section
open scoped Topology

namespace MyProjects.Topology.B

/-- A tiny "NatIso-like" record for a family of homeomorphisms natural in `Y`.
We intentionally avoid category-theoretic packaging to keep this in
`ContinuousMap`/`Homeomorph` land. -/
structure NatHomeoInY (X Z : Type*)
    [TopologicalSpace X] [TopologicalSpace Z] [LocallyCompactSpace X] where
  app : ∀ (Y : Type*) [TopologicalSpace Y] [LocallyCompactSpace Y],
    C(X × Y, Z) ≃ₜ C(X, C(Y, Z))

/-- Canonical natural family: `Y ↦ expLawCOHomeo (X, Y, Z)`. -/
def expLawCO_natIso_Y
  (X Z : Type*) [TopologicalSpace X] [TopologicalSpace Z] [LocallyCompactSpace X] :
  NatHomeoInY X Z :=
{ app := fun Y _ _ => expLawCOHomeo (X:=X) (Y:=Y) (Z:=Z) }

-- convenience simp lemmas (pointwise)

@[simp] lemma expLawCO_natIso_Y_apply
  [TopologicalSpace X] [TopologicalSpace Y] [TopologicalSpace Z]
  [LocallyCompactSpace X] [LocallyCompactSpace Y]
  (f : C(X × Y, Z)) (x : X) (y : Y) :
  (expLawCO_natIso_Y (X:=X) (Z:=Z)).app Y f x y = f (x, y) := rfl

@[simp] lemma expLawCO_natIso_Y_symm_apply
  [TopologicalSpace X] [TopologicalSpace Y] [TopologicalSpace Z]
  [LocallyCompactSpace X] [LocallyCompactSpace Y]
  (g : C(X, C(Y, Z))) (x : X) (y : Y) :
  ((expLawCO_natIso_Y (X:=X) (Z:=Z)).app Y).symm g (x, y) = g x y := rfl

/-- Y‑naturality (inner precomposition), pointwise form. -/
@[simp] lemma expLawCO_natIso_Y_precomp_right_apply
  [TopologicalSpace X] [TopologicalSpace Y] [TopologicalSpace Y'] [TopologicalSpace Z]
  [LocallyCompactSpace X] [LocallyCompactSpace Y] [LocallyCompactSpace Y']
  (τ : C(Y', Y)) (f : C(X × Y, Z)) (x : X) (y' : Y') :
  ((expLawCO_natIso_Y (X:=X) (Z:=Z)).app Y'
      (f.comp (ContinuousMap.prodMk ContinuousMap.fst (τ.comp ContinuousMap.snd)))) x y'
  =
  ((expLawCO_natIso_Y (X:=X) (Z:=Z)).app Y f) x (τ y') := by
  simpa using
    (expLawCO_precomp_right_apply (X:=X) (Y:=Y) (Y':=Y') (Z:=Z) τ f x y')

end MyProjects.Topology.B

\end{lstlisting}

\end{document}
